% ==================================================================
\section{问题重述}
% ==================================================================

\subsection{问题背景}

人在表达情绪时很少只靠一种信号：说了什么（文本）、怎么说的（语调、停顿）、说话时的表情，三者共同决定了一句话的情感色彩。多模态情感分析就是把这几路信号放在一起建模，它是具身智能、人机对话等应用绕不开的环节。

落到实际数据上，首先要解决的是“对齐”：文本以词为单位、没有时间戳，语音按毫秒分帧，视频按固定帧率抽帧，三路信号的时间尺度各不相同，必须先统一到同一套序列位置上，下游模型才能使用。在此基础上，还有两个问题特别常见。第一个是缺失：镜头被挡住几秒、背景噪声盖过了人声、语音转写漏掉了半句，都会让某一路信号在一段时间里“空掉”。很多融合模型默认三路输入都在，一旦某路局部为空，性能会掉得很厉害。第二个是可信度：模型给出“负向、强度$-1.1$”之后，使用者往往想知道它凭的是哪句话、哪一段声音或哪几帧表情，而大多数模型给不出可核对的依据。

赛题以 CMU-MOSEI 英文多模态情感数据集\cite{zadeh2018mosei}为基础，要求依次完成特征提取与时序对齐、模态局部缺失下的鲁棒预测、可解释预测三项任务，本文依次完成这三项任务。已有的经典融合方法，如张量融合网络 TFN\cite{zadeh2017tfn}、低秩融合 LMF\cite{liu2018lmf} 和跨模态 Transformer MulT\cite{tsai2019mult}，都以组织好时序的三模态特征为输入\cite{zhang2024review}，因此问题一所做的特征提取与对齐是后两问的前置环节。

\subsection{数据说明}

赛题共提供四个附件。附件1是从 CMU-MOSEI 中选出的100条原始英文视频，按原始视频编号 video\_id 分为37个子文件夹，每个片段以 clip\_id.mp4 保存，时长2.648～34.567秒；配套的标注表给出每条片段的转写文本、连续情感强度和极性标签，用于问题一。附件2是整理好的 CMU-MOSEI 特征文件，含对齐版（aligned\_50.pkl）与非对齐版两套，每套4850条样本，按训练、验证、测试划分，带有情感极性（负向、中性、正向）和连续强度（$[-3,3]$）两种标签。附件3是30条无标签的“模态局部缺失专项测试集”，缺失表现为某些模态中特征值全为零的连续区间。附件4是20条无标签的“可解释性专项测试集”，三模态完整，另附原始视频。赛题规定问题二、问题三的模型只能在附件2的训练集上学习参数、在验证集上选择结构和阈值，附件3、4只用于最终推理；分类任务用准确率和 F1 值评价，回归任务用平均绝对误差（MAE）和皮尔逊相关系数评价。

\subsection{需要解决的问题}

\textbf{问题一：}基于附件1的原始视频，建立文本、语音、视觉三模态情感特征提取与时序对齐的数学模型，明确原始样本的处理流程、各模态特征的定义与提取方法、跨模态时序对齐的规则与实现逻辑；生成的特征文件要做到样本覆盖完整、时序组织可核验、方法可复现，并以典型样本展示对齐效果。

\textbf{问题二：}针对文本、语音或视觉中部分连续时间段不可用的情形，构建鲁棒的多模态情感预测模型，在局部信息不完整时稳定输出情感极性与强度；分析缺失模态类型、缺失位置、缺失时长对预测性能的影响；用该模型对附件3完成推理。

\textbf{问题三：}针对三模态完整但决策依据难以追溯的情形，构建可解释的多模态情感预测模型，在输出极性与强度的同时，以可量化、可复核的方式给出主要参考模态、各模态作用程度，以及能对应到原始文本片段、语音时段或视觉关键帧的关键证据；用该模型对附件4完成推理。

% ==================================================================
\section{问题分析}
% ==================================================================

\subsection{问题一的分析}

问题一的难点在对齐，具体有三层。

第一层是三路信号的时间尺度和时间轴都不一样。文本是一串词，没有任何时间信息；语音以10毫秒为一帧，一条片段有几百到几千帧；视觉以0.1秒为一帧。要把它们放进同一套50个位置，首先得回答“位置代表什么”：代表一个词，还是代表一小段时间。

第二层是对齐结果要可核验。题目要求记录每个序列位置与原始素材的对应关系，这意味着每个位置都要能说清楚对应哪个词、哪段语音、哪几帧画面，不能只输出一堆对不上号的数字。

第三层是空白时段。片段里常有说话前后的静音、人转身或走出画面的时刻，把这些时段和有效内容一视同仁地铺进50个位置，会稀释有用的信息。

据此，问题一的路线是：三路特征分别用成熟的开源工具提取（BERT、librosa、YuNet），对齐部分设计一个公共的三角核重采样算子，再给出三种时间轴定义。按词位置对齐（口径H）复刻附件2的做法，便于与后两问的数据接口对接；全程50等分作为最直接的时间对齐基准；说话者人脸跟踪方案用“说话者的脸是否可见”划定有效时段，是本文针对空白时段提出的改进。最后用非零位置数、有效时长占比等指标比较三种方案，并用典型样本展示对齐效果。

\subsection{问题二的分析}

与问题一不同，问题二、三直接使用附件2提供的现成特征，重点从“特征怎么来”转到“模型怎么用”。问题二要的是“局部缺失下也能稳”。难点有三层。

第一层是训练数据里没有这种缺失。第\ref{sec:data}章会看到，附件2中的零帧几乎全是尾部填充，真正落在句子中间的连续缺失只有视觉模态零星几段；附件3却是30条全部带句中连续缺失。换句话说，题目考的失效模式在训练集里基本不存在，必须自己构造。

第二层是连续段缺失不好补。散点式的缺失可以靠相邻帧插值填回来，连续缺二三十帧就没有邻居可借了，模型得真正“推断”出洞里应有的表征。

第三层是规律分析。题目要求说清楚缺失的类型、位置、时长各自有多大影响，而附件3没有标签，只能在附件2上设计可控实验。三个因素交织在一起，需要把它们拆开单独扫描。

据此，问题二的路线是：先复现一个本身就考虑了缺失的基线模型。本文选用文献\cite{shi2026ldsd}的语言主导融合与自蒸馏模型（LDSD），它以语言为主模态、用自蒸馏重建缺失信息，思路上与题意吻合；然后用整模态缺失下的鲁棒性曲线和消融实验诊断它的短板；再围绕“连续段缺失”设计候选改进模块，用统一掩码逐个筛选；最后用三因子扫描回答缺失规律问题，并对附件3推理。

\subsection{问题三的分析}

与问题二相比，问题三的输入是完整的，难点从“预测得稳”转到了“说得清”。

“说得清”在这里有三个具体要求。一是可量化：每个模态起了多大作用，要给出数，而不是一句“主要依据文本”。二是可复核：给出的证据要能让人回到原始材料里核对。三是对得上：解释报告描述的必须就是实际做预测的那个模型，而不是另起一个近似模型。

这里有两个现实障碍。附件只提供了特征和原始视频，没有逐词的时间戳，所以“证据对应到语音时段、视觉关键帧”无法直接换算成秒；另外，注意力权重虽然现成，但已有研究\cite{jain2019attention}指出注意力分布与预测的因果贡献并不一定一致，拿它直接当解释并不稳妥。

因此问题三不改网络，而是在训练好的网络外部做“干预”：遮掉某个模态或某一段，看输出变了多少，变得越多说明越重要。证据定位则借助第\ref{sec:data}章发现的帧—词元一一对应关系，把重要的帧翻译回原文的词；语音和视觉只给出词级锚点，不去凑一个没有依据的秒数。

\subsection{求解思路}

问题一处理附件1的原始视频，问题二、三基于附件2的特征并共用同一个预测网络，整体流程如图\ref{fig:flow}所示。

\begin{figure}[htbp]
\centering
\begin{tikzpicture}[
  node distance=3.2mm,
  box/.style={draw=black!65, rounded corners=2pt, align=flush center, font=\small, inner sep=3pt, minimum height=7.5mm},
  wide/.style={box, text width=13.2cm, fill=black!4},
  q1/.style={box, text width=13.2cm, fill=green!6},
  q2/.style={box, text width=6.2cm, fill=blue!5},
  q3/.style={box, text width=6.2cm, fill=orange!7},
  head/.style={font=\small\bfseries, text width=6.2cm, align=flush center},
  arr/.style={-{Stealth[length=1.8mm]}, draw=black!70, semithick}]
\node[head, text width=13.2cm] (h1) {问题一：多模态特征提取与时序对齐（附件1）};
\node[q1, below=2mm of h1] (p1) {原始视频 $\to$ 文本 BERT（768维）、语音 librosa（74维）、视觉 YuNet（35维）};
\node[q1, below=of p1] (p2) {三角核重采样到50个位置：按词位置（口径H）\,/\,全程50等分\,/\,说话者有效时段50等分};
\node[q1, below=of p2] (p3) {对齐质量评估、典型样本验证、100条全量特征文件};
\node[wide, below=9mm of p3] (d) {数据核查（附件2、3、4）：零帧结构、有效长度、帧与词元的对应关系};
\node[wide, below=of d] (b) {基线复现：语言主导融合与自蒸馏模型 LDSD};
\node[head, below=5mm of b.south west, anchor=north west] (h2) {问题二：局部缺失下的鲁棒预测};
\node[head, below=5mm of b.south east, anchor=north east] (h3) {问题三：可解释预测};
\node[q2, below=2mm of h2] (a1) {整模态缺失下的鲁棒性与三组消融};
\node[q2, below=of a1] (a2) {候选模块 M1–M5 逐级筛选};
\node[q2, below=of a2] (a3) {LDSD-Seg（M1+M4）全量确认};
\node[q2, below=of a3] (a4) {三因子扫描：类型 × 位置 × 时长};
\node[q2, below=of a4] (a5) {附件3推理与错误归因};
\node[q3, below=2mm of h3] (c1) {L0 帧—词元回指};
\node[q3, below=of c1] (c2) {L1 干预式模态归因};
\node[q3, below=of c2] (c3) {L2 滑动窗口片段扫描};
\node[q3, below=of c3] (c4) {L3 证据呈现：词组 / 词级锚点};
\node[q3, below=of c4] (c5) {解释质量检验与附件4推理};
\draw[arr] (p1) -- (p2); \draw[arr] (p2) -- (p3);
\draw[arr, dashed] (p3) -- node[right, font=\scriptsize]{口径H与附件2的词级对齐一致（第\ref{sec:align}节验证）} (d);
\draw[arr] (d) -- (b);
\draw[arr] (b.south -| h2.north) -- (h2.north);
\draw[arr] (b.south -| h3.north) -- (h3.north);
\foreach \x/\y in {a1/a2,a2/a3,a3/a4,a4/a5,c1/c2,c2/c3,c3/c4,c4/c5} \draw[arr] (\x) -- (\y);
\end{tikzpicture}
\caption{全文求解思路（问题二、三共用基线网络，问题三直接作用于训练好的网络）}
\label{fig:flow}
\end{figure}

% ==================================================================
\section{模型假设}
% ==================================================================

\begin{enumerate}[label=(\arabic*)]
\item 附件2、3、4的特征由同一流程提取，同一序列位置上的三模态特征对应原句中的同一个词。
\item 附件3中特征全为零的连续区间就是题目所说的局部缺失；有效长度之外的零帧属于补齐填充，不算缺失。
\item 附件3、附件4的样本除缺失形态外与附件2同分布，在附件2上选定的结构和超参数可以直接用于两个专项测试集。
\item 对单个样本，遮蔽某一模态或某一片段后模型输出的变化量，可以衡量该部分对本次判断的贡献。
\item 附件2的连续强度标签是多名标注者评分的平均值，视为无偏真值；极性由强度的符号决定，0只属于中性。
\item 附件1的转写文本与语音内容一致；在缺少逐词时间戳时，词在语音时段内近似匀速分布。
\item 附件1的片段以单人讲话为主，说话者是离镜头最近、人脸面积最大的人。
\end{enumerate}

\textbf{说明。}假设(1)是后面证据回指的前提，我们没有只凭文件说明就接受它，而是在第\ref{sec:align}节用数据做了验证：语音有效长度与 BERT 词元数的相关系数为1.0000，训练、验证、测试三个划分逐条成立。

假设(4)是问题三干预式解释的基础。它常被质疑的一点是：遮蔽后的输入可能偏离训练分布，输出变化里混进了“模型没见过这种输入”的成分。对本文的模型来说这个问题要轻得多，因为训练时本来就注入了整模态和连续段缺失（第\ref{sec:q2}章），被遮蔽的输入正是模型见过的形态。

假设(5)有一个旁证：附件4中能在附件2里找回真值的5条样本，强度分别为$-0.333$、$0.667$、$0.333$、$-0.667$、$1.000$，都是$1/3$的整数倍，与“三名标注者取平均”的做法一致。

假设(6)只在问题一的方法一、二中用来给词元分配时刻。真实语速有快有慢，这个近似会带来词级的时间误差；更精确的做法是用语音识别的强制对齐得到逐词时间戳，我们把它列为改进方向。假设(7)是方法三锁定说话者的依据，对单人独白为主的 CMU-MOSEI 片段基本成立，同一画面出现两张大小相近的人脸时可能误判。

% ==================================================================
\section{符号说明}
% ==================================================================

\begin{table}[htbp]
\centering
\caption{主要符号说明}
\label{tab:symbols}
\begin{tabular}{cl}
\toprule
符号 & 含义 \\
\midrule
$m\in\{l,a,v\}$ & 模态：语言（文本）、语音、视觉 \\
$T=50$ & 对齐后的序列位置数（词级对齐） \\
$U_m^0\in\mathbb{R}^{T\times d_m}$ & 模态 $m$ 的原始特征，$d_l=768$，$d_a=74$，$d_v=35$ \\
$U_m^1$ & 注入缺失后的特征 \\
$L_m$ & 模态 $m$ 的有效长度（最后一个非零帧的位置加1） \\
$M_m\in\{0,1\}^T$ & 模态 $m$ 的缺失掩码，1 表示该帧缺失 \\
$d=128$ & 统一隐层维度（8头，每头 $d_k=16$） \\
$H_m^1$，$H_l^i$ & 模态嵌入；语言第 $i$ 个尺度的表示 \\
$H_{Hyp}^j$ & 第 $j$ 层超模态表示 \\
$E_{rec}$，$H^{rec}$ & 重建器及其输出 \\
$\theta^{tea}$，$\theta^{stu}$，$\eta$ & 教师、学生参数与 EMA 系数（$\eta=0.99$） \\
$\hat y_{reg}$，$\hat y_{cls}$ & 强度预测值与极性预测分布 \\
$\lambda$，$\gamma$ & 损失权重（$\lambda=0.5$，$\gamma=1.0$） \\
$r$ & 缺失率，含义随缺失模式而定（见第\ref{sec:missdef}节） \\
$\rho$ & 因子扫描中的缺失时长（占该模态自身内容跨度的比例） \\
$s_m$，$\mu$ & 模态作用程度；主导模态的判定裕度 \\
$\mathrm{Imp}_m(t)$，$w$ & 片段重要性；遮蔽窗口宽度 \\
\midrule
$D$，$D_{sp}$，$D_{eff}$ & 问题一：片段时长、语音时长、有效时长 \\
$L_{txt}$ & 问题一：有效词元数（含 \texttt{[CLS]}、\texttt{[SEP]}） \\
$\tau_i$，$a_k$，$b$ & 问题一：帧时间戳、对齐锚点、三角核带宽 \\
$A_s$，$\kappa$ & 问题一：说话者典型脸面积；可见判定的面积比（$\kappa=0.4$） \\
\bottomrule
\end{tabular}
\end{table}

注：此为本文的主要符号，其他符号在正文首次出现处说明。
